\documentclass[../../main.tex]{subfiles}
\graphicspath{{../../image/}{../../figure/}}
\begin{document}

\section{黏性行波分解}\label{section:ch15-黏性行波分解}

在\cref{section:ch15-黏性行波分解}中我们看到，通过把解表示为简单波的叠加 (7.8.1)，
双曲系统\eqref{align:ch15-1511}约化为弱耦合标量方程组 (7.8.6)。

以速度 $s$ 传播的黏性波是\eqref{align:ch15-1521}形如 $U(x,t)=V(x-st)$ 的解。函数 $V$ 必须满足二阶常微分方程
\begin{align}
\ddot V=[A(V)-sI]\dot V,
\label{align:ch15-1541}
\end{align}
它可改写为一阶系统
\begin{align}
\begin{cases}
\dot V=W,\\
\dot W=[A(V)-sI]W,\\
\dot s=0.
\end{cases}
\label{align:ch15-1542}
\end{align}

显然存在依赖于 $(2n+1)$ 个参数的黏性行波族，其参数为速度以及 $V$ 与 $W$ 的初值。但对当前目的而言，唯一有用的是那些 $|W|$ 在 $\mathbb{R}$ 上有界且小、$s$ 接近某特征速度 $\lambda_j^*$ 的波。它们称为\emph{黏性 $j$-波}。

对系统
\begin{align}
\begin{cases}
\dot V=W,\\
\dot W=[A^*-\lambda_j^*I]W,\\
\dot s=0,
\end{cases}
\label{align:ch15-1543}
\end{align}
它由在 $(V=U^*,W=0,s=\lambda_j^*)$ 处线性化\eqref{align:ch15-1542}得到，$W$ 有界的解轨道张成 $(n+2)$ 维超平面
\begin{align}
\mathscr{P}_j=\{(V,W,s):V\in\mathbb{R}^n,\ W=wR_j^*,\ w\in\mathbb{R},\ s\in\mathbb{R}\},
\label{align:ch15-1544}
\end{align}
它嵌入于 $\mathbb{R}^{2n+1}$ 中。由中心流形定理可知，黏性 $j$-波的轨道张成一个光滑的 $(n+2)$ 维流形 $\mathscr{M}_j$，它嵌入于 $\mathbb{R}^{2n+1}$，并在点 $(U^*,0,\lambda_j^*)$ 处与 $\mathscr{P}_j$ 相切。进而，该\emph{中心流形}可局部参数化为 $(V,w,s)$，$V\in\mathbb{R}^n$，$w\in\mathbb{R}$，$s\in\mathbb{R}$，即
\begin{align}
\mathscr{M}_j=\{(V,W,s):|V-U^*|<\varepsilon,\ W=wS_j(V,w,s),\ |w|<\varepsilon,\ |s-\lambda_j^*|<\varepsilon\},
\label{align:ch15-1545}
\end{align}
其中 $S_j$ 是光滑单位向量场，满足 $S_j(U^*,\lambda_j^*)=R_j^*$。

我们进而证明：在不变流形 $\mathscr{M}_j$ 上\eqref{align:ch15-1542}的轨道上，$w$ 满足一个微分方程。首先，由于 $|S_j|=1$，
\begin{align}
S_j^\top S_j=1,\qquad S_j^\top\dot S_j=0,\qquad S_j^\top\ddot S_j=-\dot S_j^\top\dot S_j.
\label{align:ch15-1546}
\end{align}
由于 $W=wS_j$ 满足\eqref{align:ch15-1542}$_2$，
\begin{align}
\dot wS_j+w\dot S_j=w[A-sI]S_j.
\label{align:ch15-1547}
\end{align}
将\eqref{align:ch15-1547}左端乘以 $S_j^\top$ 并利用\eqref{align:ch15-1546}，得到
\begin{align}
\dot w=(\sigma_j-s)w,
\label{align:ch15-1548}
\end{align}
其中
\begin{align}
\sigma_j(V,w,s)=S_j^\top(V,w,s)A(V)S_j(V,w,s).
\label{align:ch15-1549}
\end{align}
联合\eqref{align:ch15-1548}与\eqref{align:ch15-1547}并利用\eqref{align:ch15-1542}$_1$，
\begin{align}
[A-\sigma_jI]S_j=\dot S_j=DS_j\dot V+\dot w\partial_wS_j=w\bigl[DS_jS_j+(\sigma_j-s)\partial_wS_j\bigr].
\label{align:ch15-15410}
\end{align}
在\eqref{align:ch15-15410}中令 $w\to0$，得
\begin{align}
\sigma_j(V,0,s)=\lambda_j(V),\qquad S_j(V,0,s)=R_j(V).
\label{align:ch15-15411}
\end{align}
微分\eqref{align:ch15-1548}，
\begin{align}
\ddot w=(\sigma_jw)^\cdot-s\dot w.
\label{align:ch15-15412}
\end{align}

总之，当 $U$ 是黏性 $j$-波时，$U(x,t)=V(x-st)$ 且 $s$ 接近 $\lambda_j^*$，则
\begin{align}
\begin{cases}
\partial_xU=w_jS_j(U,w_j,s_j),\\
\partial_tU=\omega_jS_j(U,w_j,s_j),
\end{cases}
\label{align:ch15-15413}
\end{align}
其中 $s_j(x,t)=s$，$w_j(x,t)=w(x-st)$，$\omega_j(x,t)=-s_j(x,t)w_j(x,t)$。注意\eqref{align:ch15-1542}与\eqref{align:ch15-15413}联合蕴含
\begin{align}
\partial_x^2U=w_j[A(U)-s_jI]S_j(U,w_j,s_j).
\label{align:ch15-15414}
\end{align}

进而，由\eqref{align:ch15-15412}，
\begin{align}
\begin{cases}
\partial_tw_j+\partial_x(\sigma_jw_j)-\partial_x^2w_j=0,\\
\partial_t\omega_j+\partial_x(\sigma_j\omega_j)-\partial_x^2\omega_j=0.
\end{cases}
\label{align:ch15-15415}
\end{align}

现考虑\eqref{align:ch15-1521}的任一解 $U$，并把它设想为黏性波的叠加，如下所述。对上半平面任一点 $(x,t)$，寻求数 $w_j=w_j(x,t)$，$\omega_j=\omega_j(x,t)$，$s_j=s_j(x,t)$，$j=1,\cdots,n$，满足 $|w_j|<\varepsilon$，$|\omega_j|<\varepsilon$，$|s_j-\lambda_j^*|<\varepsilon$，并使得
\begin{align}
\begin{cases}
\partial_xU=\displaystyle\sum_{j=1}^nw_jS_j(U,w_j,s_j),\\[2mm]
\partial_tU=\displaystyle\sum_{j=1}^n\omega_jS_j(U,w_j,s_j).
\end{cases}
\label{align:ch15-15416}
\end{align}

事实上，假设 $|\partial_xU|$、$|\partial_tU|$ 充分小，且对任一固定的 $s_j\in\mathscr{B}_\varepsilon(\lambda_j^*)$，\eqref{align:ch15-15411}与隐函数定理联合蕴含存在唯一的 $w_j$、$\omega_j$ 使\eqref{align:ch15-15416}成立。回顾\eqref{align:ch15-15413}，上述求和中的第 $j$ 项是黏性 $j$-波的一阶逼近。接着我们考察是否能通过审慎选取 $s_j$ 达到更紧的拟合。联合\eqref{align:ch15-15416}与\eqref{align:ch15-1521}得
\begin{align}
\partial_x^2U=\sum_{j=1}^nw_j\left[A(U)+\frac{\omega_j}{w_j}I\right]S_j(U,w_j,s_j).
\label{align:ch15-15417}
\end{align}
比较\eqref{align:ch15-15417}与\eqref{align:ch15-15414}可知，若取 $s_j=-\omega_j/w_j$，则可达到二阶拟合。困难在于这可能与要求 $|s_j-\lambda_j^*|<\varepsilon$ 冲突。作为折中，我们取
\begin{align}
s_j=\lambda_j^*-\theta\left(\lambda_j^*+\frac{\omega_j}{w_j}\right),
\label{align:ch15-15418}
\end{align}
其中 $\theta$ 是光滑“截断”函数，满足
\begin{align}
\theta(r)=
\begin{cases}
r&|r|\leqslant\delta_1,\\
0&|r|>3\delta_1,
\end{cases}
\qquad|\theta'(r)|\leqslant1,\qquad|\theta''(r)|\leqslant\frac4{\delta_1},
\label{align:ch15-15419}
\end{align}
其中 $\delta_1$ 为某个小正数。于是当 $-\omega_j/w_j$ 取接近 $\lambda_j^*$ 的值时，$s_j=-\omega_j/w_j$；反之，当 $-\omega_j/w_j$ 远离 $\lambda_j^*$ 时，$s_j$ 取常值 $\lambda_j^*$。

经过冗长分析（它依赖于函数 $S_j$ 的性质，见\cref{section:ch15-注记}所引文献）可以证明：只要 $|U-U^*|$、$|\partial_xU|$、$|\partial_x^2U|$、从而 $|\partial_tU|$ 对每个 $(x,t)$ 都充分小，就存在唯一的一组 $(w_j,\omega_j)$，$j=1,\cdots,n$，满足\eqref{align:ch15-15416}与\eqref{align:ch15-15418}。进而，函数 $(w_j(x,t),\omega_j(x,t))$，$j=1,\cdots,n$，是 $C^{1,1}$ 光滑的。此外，对照\cref{lemma:ch15-扩散效应估计}的设定与记号，当\eqref{align:ch15-1523}在区间 $(0,T)$ 上成立、$\bar t<T\leqslant\infty$（此时\eqref{align:ch15-1524}对任意 $t\in[\bar t,T)$ 成立）时，
\begin{align}
\begin{cases}
\displaystyle\sum_{j=1}^n\left\{\|w_j(\cdot,t)\|_{L^1(-\infty,\infty)}+\|\omega_j(\cdot,t)\|_{L^1(-\infty,\infty)}\right\}=O(\delta_0),\\[2mm]
\displaystyle\sum_{j=1}^n\left\{\|w_j(\cdot,t)\|_{L^\infty(-\infty,\infty)}+\|\omega_j(\cdot,t)\|_{L^\infty(-\infty,\infty)}\right\}=O(\delta_0^2),\\[2mm]
\displaystyle\sum_{j=1}^n\left\{\|\partial_xw_j(\cdot,t)\|_{L^1(-\infty,\infty)}+\|\partial_x\omega_j(\cdot,t)\|_{L^1(-\infty,\infty)}\right\}=O(\delta_0^2),\\[2mm]
\displaystyle\sum_{j=1}^n\left\{\|\partial_xw_j(\cdot,t)\|_{L^\infty(-\infty,\infty)}+\|\partial_x\omega_j(\cdot,t)\|_{L^\infty(-\infty,\infty)}\right\}=O(\delta_0^3),
\end{cases}
\label{align:ch15-15420}
\end{align}
对 $t\in[\bar t,T)$ 一致成立。

如上所见，当 $U$ 恰为一个黏性 $j$-波时，$w_j$ 与 $\omega_j$ 满足\eqref{align:ch15-15415}。对一般解 $U$，取而代之的是带源项 $\varphi_j$ 与 $\psi_j$ 的方程\eqref{align:ch15-1526}。我们期望用黏性波对 $U$ 的逼近足够紧，以使 $\varphi_j$ 与 $\psi_j$ “小”。

经过冗长计算（见\cref{section:ch15-注记}所引文献）可得
\begin{align}
(\varphi_j,\psi_j)&=O(1)\sum_{i\neq k}\bigl(|w_iw_k|+|\omega_i\omega_k|+|w_i\omega_k|+|w_i\partial_xw_k|+|w_i\partial_x\omega_k|+|\omega_i\partial_xw_k|\bigr)\nonumber\\
&\quad+O(1)\sum_i|\omega_i\partial_xw_i-w_i\partial_x\omega_i|\nonumber\\
&\quad+O(1)\sum_i\left|w_i\partial_x\left(\frac{\omega_i}{w_i}\right)\right|^2\chi_{\{|\lambda_i^*+\omega_i/w_i|<3\delta_1\}}\nonumber\\
&\quad+O(1)\sum_i\bigl(|\partial_xw_i|+|\partial_x\omega_i|\bigr)|\omega_i+s_iw_i|.
\label{align:ch15-15421}
\end{align}

\eqref{align:ch15-15421}右端四项估计了\eqref{align:ch15-1526}相对于单个黏性 $j$-波情形\eqref{align:ch15-15415}的“偏差”，其来源如下：
\begin{enumerate}
\item\label{item:ch15-154-a} 第一项解释横波相互作用：属于不同特征族、因而以不同速度传播的黏性波相互作用，对 $\varphi_j$ 与 $\psi_j$ 产生二次阶贡献。
\item\label{item:ch15-154-b} 第二、三项解释同族波之间的相互作用：在两点 $x$ 与 $y$ 处逼近剖面 $U(\cdot,t)$ 的黏性 $i$-波以不同速度 $s_i(x,t)$、$s_i(y,t)$ 传播，因而可能相互作用。估计中的关键因子是 $\partial_xs_i$，它刻画了 $s_i$ 的变化率。
\item\label{item:ch15-154-c} 第四项解释通过\eqref{align:ch15-15418}而非 $-\omega_i/w_i$（对黏性 $i$-波本应如此）选取 $s_i$ 所造成的“误差”。注意当 $s_i=-\omega_i/w_i$ 时此项消失。
\end{enumerate}

以下三节将估计\eqref{align:ch15-15421}右端。目标是验证\cref{section:ch15-证明路线图}中的断言，即若\eqref{align:ch15-1523}在 $(0,T)$ 上成立，则\eqref{align:ch15-1527}蕴含\eqref{align:ch15-1528}。

\end{document}
